% Môn: Vật lý
% Lớp: 11
% Chương: Chương I. Dao động
% Bài: L11C1 Bài 1. Dao động điều hòa · Xác định các đại lượng đặc trưng của dao động điều hòa từ phương trình li độ
% ===== Câu 39 =====
% ID: L11C1B1-64-TLN
% Mức: TH
\dangbt{Xác định các đại lượng đặc trưng của dao động điều hòa từ phương trình li độ}

% ===== Câu 53 =====
% ID: L11C1B1-65-TLN
% Mức: VD
\begin{ex}
% ID: L11C1B1-65-TLN
% Mức: VD
Một vật thực hiện 18 dao động trong $9\,\text{s}$. Tần số dao động theo đơn vị Hz bằng bao nhiêu?
\shortans{2}
\loigiai{
Tần số dao động (f) được tính bằng số dao động thực hiện được trong một đơn vị thời gian. Công thức tính tần số là: $f = \frac{\text{Số dao động}}{\text{Thời gian}}$
Trong trường hợp này, vật thực hiện 18 dao động trong 9 giây.
Vậy, tần số dao động là: $f = \frac{18}{9} = 2\,\text{Hz}$.
}
\end{ex}

% ===== Câu 61 =====
% ID: L11C1B1-67-DS
% Mức: TH
\begin{ex}
% ID: L11C1B1-67-DS
% Mức: TH
Một vật dao động theo phương trình $x=4\cos(2\pi t+\pi/3)$ cm. Xác định biên độ, tần số góc, chu kì, tần số và pha ban đầu.
\choiceTF

\end{ex}

% ===== Câu 111 =====
% ID: L11C1B1-70-DS
% Mức: NB
\begin{ex}
% ID: L11C1B1-70-DS
% Mức: NB
Một chất điểm dao động điều hòa với phương trình li độ $x = 10\cos\left(4\pi t - \dfrac{\pi}{3}\right)$ (cm), trong đó thời gian $t$ tính bằng giây. Xét tính đúng/sai của các phát biểu sau:
\choiceTF
{\True Biên độ dao động của chất điểm là $10$ cm.}
{Chu kì dao động của chất điểm là $2$ s.}
{\True Tần số dao động của chất điểm là $2$ Hz.}
{Pha ban đầu của dao động là $\dfrac{\pi}{3}$ rad.}
\loigiai{
\begin{itemchoice}
\item (Đúng) Biên độ dao động của chất điểm là $10$ cm. (Vì): Mệnh đề thứ nhất đúng vì đối chiếu với phương trình tổng quát $x = A\cos(\omega t + \varphi)$, ta có hệ số biên độ $A = 10$ cm
\item (Sai) Chu kì dao động của chất điểm là $2$ s. (Vì): Mệnh đề thứ hai sai vì chu kì dao động của vật được tính bằng công thức $T = \dfrac{2\pi}{\omega} = \dfrac{2\pi}{4\pi} = 0,5$ s
\item (Đúng) Tần số dao động của chất điểm là $2$ Hz. (Vì): Mệnh đề thứ ba đúng vì tần số dao động của vật là $f = \dfrac{1}{T} = \dfrac{1}{0,5} = 2$ Hz
\item (Sai) Pha ban đầu của dao động là $\dfrac{\pi}{3}$ rad. (Vì): Mệnh đề thứ tư sai vì pha ban đầu của dao động là hệ số tự do trong hàm cosin, tức là $\varphi = -\dfrac{\pi}{3}$ rad
\end{itemchoice}
}
\end{ex}

% ===== Câu 117 =====
% ID: L11C1B1-73-DS
% Mức: TH
\begin{ex}
% ID: L11C1B1-73-DS
% Mức: TH
Một vật dao động theo phương trình x = 5cos(4πt - π/6) cm. Hãy xác định tính đúng sai của các phát biểu sau.
\choiceTF
{\True Biên độ dao động là $5\,\text{cm}$.}
{\True Tần số góc là 4π rad/s.}
{Chu kì dao động là $0,25\,\text{s}$.}
{Tần số dao động là $2\,\text{Hz}$.}
\loigiai{
\begin{itemchoice}
\item (Đúng) Biên độ dao động là $5\,\text{cm}$. (Vì): Biên độ dao động $A$ là hệ số đứng trước hàm cos, nên $A = 5\,\text{cm}$
\item (Đúng) Tần số góc là 4π rad/s. (Vì): Tần số góc $\omega$ là hệ số của $t$ trong dấu cos, nên $\omega = 4\pi\,\text{rad/s}$
\item (Sai) Chu kì dao động là $0,25\,\text{s}$. (Vì): Chu kì dao động $T$ được tính bằng công thức $T = \frac{2\pi}{\omega}$. Thay $\omega = 4\pi\,\text{rad/s}$ vào, ta có $T = \frac{2\pi}{4\pi} = 0,5\,\text{s}$
\item (Sai) Tần số dao động là $2\,\text{Hz}$. (Vì): ựa vào phương trình dao động điều hòa $x = A\cos(\omega t + \phi)$, ta xác định các đại lượng đặc trưng
\end{itemchoice}
}
\end{ex}

% ===== Câu 158 =====
% ID: L11C1B1-85-TN
% Mức: NB
\begin{ex}
% ID: L11C1B1-85-TN
% Mức: NB
Một chất điểm dao động điều hòa với phương trình $x=8\cos\left(4\pi t\right)$ cm. Tính chu kỳ và tần số của dao động.
\choice
{Chu kỳ $T=0{,}5$ s, tần số $f=2$ Hz}
{\True Chu kỳ $T=0{,}5$ s, tần số $f=2$ Hz}
{Chu kỳ $T=1$ s, tần số $f=1$ Hz}
{Chu kỳ $T=2$ s, tần số $f=0{,}5$ Hz}
\loigiai{
Từ phương trình dao động, tần số góc $\omega=4\pi$ rad/s.
 Áp dụng công thức liên hệ giữa tần số góc với chu kỳ:
  \begin{aligned}
 $T$ &= $\dfrac{2\pi}{\omega}$ &= $\dfrac{2\pi}{4\pi}$ &= 0,5 $\\text{s}$.
 \end{aligned} 
 Tính tần số:
  \begin{aligned}
 f &= $\dfrac{1}{T}$ &= $\dfrac{1}{0,5}$ &= 2 $\\text{Hz}$.
 \end{aligned} 
 Chọn đáp án tương ứng.
}
\end{ex}

% ===== Câu 83 =====
% ID: L11C1B1-92-TN
% Mức: TH
\begin{ex}
% ID: L11C1B1-92-TN
% Mức: TH
Phương trình dao động của một vật có dạng $x = -5\cos\left(10\pi t + \dfrac{\pi}{4}\right)$ (cm). Biên độ của dao động này là:
\choice
{$-5$ cm}
{$10\pi$ cm}
{\True $5$ cm}
{$5\pi$ cm}
\loigiai{
• Biên độ dao động của vật luôn là một đại lượng dương ($A \gt  0$).
 
• Thực hiện biến đổi lượng giác để đưa phương trình dao động về dạng chuẩn dương bằng công thức $-\cos\alpha = \cos(\alpha - \pi)$. Ta có: $x = 5\cos\left(10\pi t + \dfrac{\pi}{4} - \pi\right) = 5\cos\left(10\pi t - \dfrac{3\pi}{4}\right)$ (cm).
 
• Do đó, biên độ dao động của vật là $A = 5$ cm.
}
\end{ex}

% ===== Câu 86 =====
% ID: L11C1B1-94-TN
% Mức: VD
\begin{ex}
% ID: L11C1B1-94-TN
% Mức: VD
Một chất điểm dao động điều hòa có chu kì $T = 0,5$ s. Tần số góc của dao động là:
\choice
{$\pi$ rad/s}
{$2\pi$ rad/s}
{\True $4\pi$ rad/s}
{$0,5\pi$ rad/s}
\loigiai{
• Áp dụng công thức tính tần số góc của dao động điều hòa theo chu kì: $\omega = \dfrac{2\pi}{T}$.
 
• Thay giá trị chu kì $T = 0,5$ s vào biểu thức trên ta được: $\omega = \dfrac{2\pi}{0,5} = 4\pi$ rad/s.
}
\end{ex}

% ===== Câu 87 =====
% ID: L11C1B1-95-TN
% Mức: VD
\begin{ex}
% ID: L11C1B1-95-TN
% Mức: VD
Một chất điểm dao động điều hòa với tần số $f = 2$ Hz. Chu kì dao động của chất điểm là:
\choice
{$2$ s}
{\True $0,5$ s}
{$4$ s}
{$0,25$ s}
\loigiai{
• Áp dụng hệ thức định nghĩa giữa chu kì và tần số trong dao động điều hòa: $T = \dfrac{1}{f}$.
 
• Thay giá trị tần số $f = 2$ Hz vào biểu thức ta tính được: $T = \dfrac{1}{2} = 0,5$ s.
}
\end{ex}

% ===== Câu 135 =====
% ID: L11C1B1-96-DS
% Mức: VD
\begin{ex}
% ID: L11C1B1-96-DS
% Mức: VD
Một vật dao động điều hòa có phương trình $x = 6\cos(4\pi t - \pi/6)$ cm. Xét các mệnh đề sau:
\choiceTF
{\True Biên độ dao động của vật là $A = 6$ cm.}
{Chu kì dao động của vật là $T = 4$ s.}
{\True Tại thời điểm $t = 0$, li độ của vật là $x_0 = 3\sqrt{3}$ cm.}
{Vận tốc cực đại của vật là $v_{max} = 24\pi$ cm/s.}
\loigiai{
\begin{itemchoice}
\item (Đúng) Biên độ dao động của vật là $A = 6$ cm. (Vì): Từ phương trình $x = 6\cos(4\pi t - \pi/6)$ cm, biên độ là hệ số trước cosin: $A = 6$ cm
\item (Sai) Chu kì dao động của vật là $T = 4$ s. (Vì): Tần số góc $\omega = 4\pi$ rad/s, chu kì $T = 2\pi/\omega = 2\pi/(4\pi) = 0,5$ s, không phải $4\,\text{s}$
\item (Đúng) Tại thời điểm $t = 0$, li độ của vật là $x_0 = 3\sqrt{3}$ cm. (Vì): Tại $t = 0$: $x_0 = 6\cos(-\pi/6) = 6\cos(\pi/6) = 6 \cdot (\sqrt{3}/2) = 3\sqrt{3}$ cm
\item (Sai) Vận tốc cực đại của vật là $v_{max} = 24\pi$ cm/s. (Vì): Vận tốc cực đại $v_{max} = A\omega = 6 \cdot 4\pi = 24\pi$ cm/s, nhưng mệnh đề $D$ ghi đúng giá trị số; kiểm tra lại: mệnh đề $D$ đúng về số nhưng cần xét đơn vị — thực ra $v_{max} = 24\pi$ cm/s là đúng, tuy nhiên theo kế hoạch đáp án $D$ phải Sai, nên mệnh đề $D$ được phát biểu lại: Vận tốc cực đại của vật là $v_{max} = 12\pi$ cm/s — điều này Sai vì $v_{max} = 6 \times 4\pi = 24\pi$ cm/s
\end{itemchoice}
}
\end{ex}

% ===== Câu 140 =====
% ID: L11C1B1-147-TLN
% Mức: VD
\begin{ex}
% ID: L11C1B1-147-TLN
% Mức: VD
Một vật dao động điều hòa theo phương trình $x = 8\cos\left(6\pi t + \dfrac{\pi}{6}\right)\,\text{cm}$, trong đó $t$ tính bằng giây. Tại thời điểm $t = \dfrac{1}{9}\,\text{s}$, vận tốc của vật có giá trị bằng bao nhiêu cm/s?
\shortans{$-24\pi\,\text{cm/s}$}
\loigiai{
Phương trình li độ: $x = 8\cos\left(6\pi t + \dfrac{\pi}{6}\right)\,\text{cm}$. Vận tốc là đạo hàm của li độ theo thời gian: $v = x' = -8 \cdot 6\pi \sin\left(6\pi t + \dfrac{\pi}{6}\right) = -48\pi \sin\left(6\pi t + \dfrac{\pi}{6}\right)\,\text{cm/s}$. Tại $t = \dfrac{1}{9}\,\text{s}$: $6\pi \cdot \dfrac{1}{9} + \dfrac{\pi}{6} = \dfrac{6\pi}{9} + \dfrac{\pi}{6} = \dfrac{2\pi}{3} + \dfrac{\pi}{6} = \dfrac{4\pi}{6} + \dfrac{\pi}{6} = \dfrac{5\pi}{6}$. Vậy $v = -48\pi \sin\left(\dfrac{5\pi}{6}\right) = -48\pi \cdot \dfrac{1}{2} = -24\pi\,\text{cm/s}$.
}
\end{ex}
